\chapter{Parte 1}

\section{Exercício 1}

Considere uma lei de conservação genérica dada por

\[ \pd{u}{t} + \pd{f}{x} = 0 \]

Aproximamos $u$ e $f$ por $u_h$ e $f_h$, o que resulta em um Resíduo $R$ dado por

\[ R = \pd{u_h}{t} + \pd{f_h}{x} \]

Usando a técnica dos resíduos ponderados, temos que em cada elemento $k$

\[ \int_{D^k} R \cdot v \, dx = 0 \]

Fazendo $v$ como a função de Lagrange e substituindo a expressão de $R$, temos 

\[ \int_{D^k} \left( \pd{u_h}{t} + \pd{f_h}{x} \right)  L_j \, dx = 0 \]

\[ \int_{D^k} \pd{u_h}{t} L_j  \, dx + \int_{D^k} \pd{f_h}{x} L_j \, dx = 0 \]

Aplica integral por partes na segunda parcela da soma

\[ \int_{D^k} \pd{u_h}{t} L_j  \, dx = \int_{D^k} \od{L_j}{x} f \, dx - \left[ L_j f \right]_{x_L}^{x_R}  = 0 \]

\textcolor{red}{Justificar segunda integral por partes, completar}

\section{Exercício 2}

Dentro de um elemento $D^k$, aproximamos $u(x,t)$ por uma interpolação 
por polinômios de Lagrange (nodal)

\[ u(x,t) = u_h(x,t) = \sum_{i=1}^{N_p} u_i(t) L_i(x) \]

Logo, 

\[ \pd{u_h}{t} = \sum_{i=1}^{N_p} \od{u_i}{t} L_i \]

Substitui isso na forma fraca

\[ \int_{D^k} \sum_{i=1}^{N_p} \od{u_i}{t} L_i L_j  \, dx = \int_{D^k} \od{L_j}{x} f \, dx - \left[ L_j f \right]_{x_L}^{x_R}  = 0 \]

\[ \sum_{i=1}^{N_p} \od{u_i}{t} \int_{D^k} L_i L_j  \, dx = \int_{D^k} \od{L_j}{x} f \, dx - \left[ L_j f \right]_{x_L}^{x_R}  = 0 \]

Definindo

\[ M_{ij} = \int_{D^k} L_i L_j  \, dx \]

Temos

\[ M \od{u}{t} = \dots \]

Agora analisamos a segunda integral da forma fraca. Substituindo $f = au$, temos 

\[ \int_{D^k} \od{L_j}{x} a u \, dx  \]

Substituindo $u_h$, 

\[ \int_{D^k} \od{L_j}{x} a \sum_{i=1}^{N_p} u_i(t) L_i(x) \, dx  \]

\[ a  \sum_{i=1}^{N_p} u_i(t) \int_{D^k} \od{L_j}{x}  L_i(x) \, dx  \]

Definindo

\[ S_{ij} = \int_{D^k} \od{L_j}{x}  L_i(x) \, dx \]

Temos

\[  M \od{\mathbf{u}}{t} = a S \mathbf{u} - \mathbf{F} \]

Usar a representação nodal é mais conveniente do ponto de vista de código 
uma vez que armazenamos os pontos (nós) em que a função $u_h$ é avaliada.

\section{Exercício 3}

Quando $\alpha = 1$, a segunda parcela se anula e temos apenas a média aritmética das bordas
de ambos
elementos vizinhos.

Quando $\alpha = 0$, o fluxo leva em consideração a direção de propagação por meio 
dos termos $\hat{n}^+$ e $\hat{n}^-$.

\section{Exercício 4}

Os pontos GLL são escolhidos ao invés dos pontos equidistantes pois 
ele apresentam melhor desempenho no método numérico. 
Além disso, a matriz M se torna aproximadamente diagonal ao escolher
os pontos GLL, o que melhora o desempenho computacional.

Os pontos GLL são tabelados, tornando fácil seu uso.

\section{Exercício 5}

O método RK4 é alta ordem, sendo natural seu uso com o DGTD que também é de alta
ordem. Usar o método de Euler após o obter a forma semi-discreta iria jogar fora 
todo o ganho de acurácia que obtemos com o DGTD.

Na forma semi-discreta, temos a forma 

\[  M \od{\mathbf{u}}{t} = a S \mathbf{u} - \mathbf{F} \]

\[  \od{\mathbf{u}}{t} = (M)^{-1} a S \mathbf{u} - \mathbf{F} \]

\[  \od{\mathbf{u}}{t} = R(\mathbf{u}) \]

\noindent que pode ser resolvido através de 

\[
\frac{d\mathbf{u}}{dt} = R(\mathbf{u})
\]

\[
\mathbf{k}_1 = R(\mathbf{u}^n)
\]

\[
\mathbf{k}_2 = R\left(\mathbf{u}^n + \frac{\Delta t}{2}\mathbf{k}_1\right)
\]

\[
\mathbf{k}_3 = R\left(\mathbf{u}^n + \frac{\Delta t}{2}\mathbf{k}_2\right)
\]

\[
\mathbf{k}_4 = R\left(\mathbf{u}^n + \Delta t\,\mathbf{k}_3\right)
\]

\[
\mathbf{u}^{n+1}
=
\mathbf{u}^n
+
\frac{\Delta t}{6}
\left(
\mathbf{k}_1
+
2\mathbf{k}_2
+
2\mathbf{k}_3
+
\mathbf{k}_4
\right)
\]

\textcolor{red}{Mostra pq usa low storage RK4 em vez do RK4 convencional, qual
a diferença}

\section{Exercício 6}